\section{什么是微分？}

\noindent\textbf{1. 线性化}\quad 现在进一步分析上一节的两个例子。

在光的折射问题中，我们发现光走过折线 $APB$（图3-1-3）所需的时间 $f(x)$（上节(1)式）适合以下关系式
\[
 f(x+dx)-f(x)=\left[\frac{\cos\theta}{v_1}-\frac{\cos\varphi}{v_2}\right]dx+\text{高阶无穷小量}.
\]
对于行星绕日问题，若记动径扫过的扇形面积为 $S(t)$（假设 $t=0$ 时刻径位置在 $OA$，参看图3-1-1），则在有引力存在情况下，在 $dt$ 时间内，扫过的面积应是 $S(t+dt)-S(t)$，而且除了一个高阶无穷小以外，这个差应为扇形 $OBD$（$OB$ 为动径在时刻 $t$ 的位置），其顶角为 $d\varphi=\angle BOD$，半径为 $r$，因此面积为 $\frac12r^2d\varphi=\frac12r^2\dot\varphi dt$，于是我们又有
\[
 S(t+dt)-S(t)=\frac12r^2\dot\varphi dt+\text{高阶无穷小量}.
\]
在这两个例子中，我们都是把注意力放在右方第一项上，在折射问题中，由费马原理得出
\[
 \frac{\cos\theta}{v_1}-\frac{\cos\varphi}{v_2}=0,
\]
由此即得关于折射的斯涅尔定律；在行星运动问题中，我们得到了瞬时的面积速度
\[
 \frac{dS(t)}{dt}=\frac12r^2\dot\varphi=\text{常数}.
\]

概括这两个例子所用的数学方法，则都是对可求导的函数 $f(x)$ 使用了以下的公式
\begin{equation}
 f(x+h)-f(x)=f'(x)h+o(h).
 \tag{1}
\end{equation}
在所有关于微积分的教本中都称右方第一项为 $f(x)$ 在 $x$ 点的微分：
\begin{equation}
 df(x)=f'(x)h=f'(x)dx.
 \tag{2}
\end{equation}
然后就是一连串无休止的争论：$df$ 是不是无穷小量，$dx$ 是不是无穷小量？无论如何，把 $h$ 写成 $dx$ 总似乎有某种暗示：对于相信 $h$ 是无穷小的人，把它写成 $dx$ 似乎是暗示了它确实是无穷小量，否则从哪里谈得起“高阶无穷小”？对于为 $dx$ 是无穷小所带来的神秘气氛而迷惑的人来说，把 $dx$ 写成 $h$ 又似乎是在暗示 $dx$ 其实也就是普通的量，而可以如处理通常的数一样去对待它。特别是要注意 $h\ne0$（否则例如在行星问题中，动径被永远地冻结在时刻 $t$ 的位置不动了，因此也就无所谓“扫过一个扇形”了），所以可以用 $h$ 作分母去作除法了。而真正给出“致命一击”的仍然是伯克莱大主教：你们不是一再以数学的严格性而自傲吗？牛顿自己就说“在数学中最微小的误差也不可以忽略”，那么为什么又把高阶无穷小量丢掉了呢？费马先是把 $h$ 写成 $E$，还说明了 $E\ne0$，可是 $E^2$ 却被他丢掉了。这些问题困惑着三百年前的牛顿直到今天的大学生。

然而由牛顿的时代到19世纪晚期数学的进步都表明了，上面说的问题根本不是问题。读者可能以为这是言过其实，我们不妨读一下魏尔斯特拉斯1861年在德国柏林皇家工科大学讲课的笔记。下面是他的学生，著名数学家施瓦茨（H.~A.~Schwarz，我们常用的施瓦茨不等式就出自他手）的记录，引文录自李文林编《数学珍宝》，科学出版社，1998年，684--685页，凡下加波线的都是魏尔斯特拉斯的原话，其余是本书作者的说明：

\begin{quote}
“当 $x$ 变至 $x+h$ 时函数 $f(x)$ 所产生的全改变量 $f(x+h)-f(x)$ 一般（当 $f$ 可求导时就行——本书作者注）可分解为两个部分，第一部分正比于自变量的改变量 $h$，因此它由 $h$ 和不依赖于 $h$ 的一个因子（即关于 $h$ 为常量）（按：这里指(1)中的 $f'(x)$——本书作者注）组成，从而它当 $h$ 变为无穷小时也变为无穷小（魏尔斯特拉斯这里并没有把无穷小当成一个“怪物”，更不讨论它是否“不可分量”等等。在这个笔记前一点，他就对“变为无穷小”下了定义：“如果一个量的绝对值能变得小于任意选定的无论怎样小的量，则我们说它能变为无穷小”，例如 $|x|<\varepsilon$，$\varepsilon-\delta$ 的记号就是他的创造。关于 $\varepsilon-\delta$ 第一次清楚的表述也见于这份讲课的记录。见同书683页——本书作者注），或随 $h$ 同时变为无穷小。然而，另一部分当 $h$ 变为无穷小时不仅自身变为无穷小（即当除以 $h$ 后仍变为无穷小——本书作者注）。

……

函数全改变量的正比于自变量的改变量的第一部分，称为微分改变量或微分，并以刻画其特征的 $d$ 置于函数前表示，而前缀 $\Delta$ 则表示全改变量（这说明 $df$ 与 $\Delta f$ 只是两个符号，使用了这个记号后，(1)将表示为
\[
 \Delta f=df+o(h).
\]
这里没有任何一点神秘色彩，也没有丢掉任何一个不为0的量——本书作者注）。对 $h$ 我们也类似地写为 $dx$，因为 $x$ 的最简单的函数是 $x$ 自身，而 $dx$ 是完全不依赖于 $x$ 的能变为无穷小的量（就是说，若令 $f(x)\equiv x$，这是 $x$ 的最简单的函数，对于这个函数 $o(h)\equiv0$。其实，要想 $o(h)\equiv0$，必要充分条件是 $f(x)=Ax+C$，$A$ 和 $C$ 是任意常数。对于这个函数 $\Delta f=f(x+h)-f(x)=A(x+h)-Ax=Ah$。另一方面，因 $o(h)\equiv0$，故由(3)又有 $\Delta x=dx$。总之，对于自变量 $x$ 而言，$\Delta x=dx=h$，这样就有了(2)式，而且，不存在 $dx$ 是不是无穷小量的问题，它只
\end{quote}

\noindent 是能变为无穷小的量，即可以趋于零的量。十分重要的是，$dx$ 完全不依赖于 $x$，这一点我们在正文中要详细地讲解——本书作者注）。$dx$ 变 $h$ 取得越小，微分改变量与全改变量的差异越少，通过减小 $dx$ 能使此差别小于每个无论怎样小的量，因此人们定义微分为函数当其自变量改变一无穷小量时所产生的改变量。”

这一段话把古典的，即通常微积分教科书中所讲的微分究竟是什么说得再清楚不过了。概括起来就是：

1. $\Delta f$ 可分为两部分，一部分与 $dx=h$ 成线性关系的称微分，记作 $df$；另一部分对 $h$ 是高阶无穷小量，这里什么都没有舍去。所以 $df$ 是 $\Delta f$ 的线性部分。许多书上都爱说“主要部分”，这至少是不准确的。因为例如表达式 $3h+h^2$ 比之 $h$，自然 $3h$ 是主要部分，因为 $3h/(3h+h^2)\to1\,(h\to0)$，而另一部分 $h^2$ 则有 $h^2/(3h+h^2)\to0\,(h\to0)$。但是它也可写为 $\left(3h+\frac{h^2}{2}\right)+\frac{h^2}{2}$，仍有 $\left(3h+\frac{h^2}{2}\right)/(3h+h^2)\to1$，$\frac{h^2}{2}/(3h+h^2)\to0$。那么为什么不说 $3h+\frac{h^2}{2}$ 是主要部分呢？再说，若 $f'$ 在 $x$ 点为0，则由(1)式在此点 $f(x+h)-f(x)=o(h)$，而我们不能说0是 $\Delta f$ 的主要部分。线性部分这种说法是大有深意的。

2. $\Delta f,df,dx$ 无非是记号，它们都是普通的量，使用这些为了刻画它们是怎样来的，这里面没有任何神秘的东西。

3. 有些时候我们要令 $h$ 很小，这时可以看到 $df$ 也很小，误差项 $o(h)$ 对某个 $h$，哪怕是很小的 $h$，都不一定在数值上更小，但一般说来会是更小的。在有必要比较 $f(x+h)-f(x)$ 与 $f'(x)h$ 时，确实需要让 $h$ 取很小的值。再说 $f(x)$ 可能只定义在某个区间 $[a,b]$ 上，若 $h$ 不是很小，$x+h$ 就会越出这个区间。特别是，若 $x=b$，必需令 $h\le0$，否则 $f(x+h)=f(b+h)$ 根本没有定义。此外，在需要比较 $f(x+h)-f(x)$ 与 $f'(x)h$ 作为无穷小量的阶时才有需要使 $h=dx$ 为无穷小量。正因为在实际使用起微分概念时，时常要让 $h$ 很小或者甚至 $h\to0$，这样一种习惯的力量使得“人们定义微分为函数当其自变量改变一无穷小量时所产生的改变量”，而作为一个数学概念 $df$ 就是 $\Delta f$ 的线性部分，即“正比于自变量的改变量 $h$”的那一部分：$f'(x)h$。这里对 $h$ 之大小没有任何限制。

然而一切数学概念凡十分重要的，如微分等等，都一定是随历史发展的。微分的现代的定义与上述古典定义之差别在于，现在把微分看作一个映射。于是若有一个定义在区间 $(a,b)$ 上的函数 $f(x)$，或者说有一个映射 $f:(a,b)\to\mathbb R$，如果 $f(x)$ 在 $(a,b)$ 上可求导，我们就给出

\noindent\textbf{定义1}\quad 对上述 $f(x)$ 以及 $x\in(a,b)$，若有一线性映射 $A:\mathbb R\to\mathbb R$，$h\mapsto Ah$，使得
\begin{equation}
 f(x+h)-f(x)=Ah+o(h),
 \tag{4}
\end{equation}
则映射 $A$ 称为 $f$ 在 $x$ 点的微分。

由函数导数之定义即知，若令 $A=f'(x)$，则这个 $A$ 就是微分。这里注意几点：

1. 现在微分定义为一个算子、一个运算或称为一个映射。现在，$f'(x)$ 并不表示一个数，而是表示乘以 $f'(x)$ 这个数，因此是一个映射。这个映射称为原映射 $f$ 的切映射。古典定义把切映射作用于 $h\in\mathbb R$ 之结果称为微分，现在则把切映射本身称为微分。因此古典的微分 $df=Ah=f'(x)h$，现在的微分则是 $df=f'(x)$。（我们时常还是保留 $df$ 这个记号，这又会造成相当多的麻烦）。二者的区别有如函数值与函数关系。

2. 把 $f$ 也看成映射，于是(4)就表示，用切映射 $f'(x)$ 去代替原来的映射 $f$，这个过程就称为线性化。这是数学物理一个根本的处理问题的方法。

3. 在以上的讨论中我们必须把 $x$ 与 $h$ 分开。魏尔斯特拉斯的那一段话中就提到 $dx=h$ 是完全不依赖于 $x$ 的。现在从几何上来看这个问题。$y=f(x)$ 是一条曲线，它定义在 $\mathbb R^1$——$x$ 轴——上的开区间 $(a,b)$ 上，但是它在 $P$ 点（自变量为 $x$ 处）有一切线，这切线可以无限延伸，其定义域是 $-\infty<h<+\infty$，即图3-2-1上的斜线，不过我们把它的原点 $h=0$ 移到了 $x$ 处。曲线与切线都是映射。前者是定义在 $(a,b)$ 上的非线性映射，后者是定义在 $-\infty<h<+\infty$ 上的线性映射。我们在图上画出了它们的图像（graph）。这两个映射之差当 $h\to0$ 时是 $h$ 的高阶无穷小。曲线的定义域 $(a,b)$ 是 $\mathbb R^1$——$x$ 轴——的一部分开集，切线的定义域则是整个 $\mathbb R^1$ 空间 $-\infty<h<+\infty$，我们称此空间为曲线在 $P$ 点的切空间，记作 $T_P$。我们对上面用的图像二字还要作一些解释，若有一个将 $E$ 的子集 $A$（称为定义域）映入 $F$ 内的映射 $\Phi:A\subset E\to F$，所谓图像就是乘积集合 $E\times F$ 的子集 $\operatorname{graph}(\Phi)=\{(x,y):x\in A,\ y=\Phi(x)\}$。所以上图中有两个映射，一是 $f:(a,b)\to\{y\text{轴}\}$，其自变量用 $x$ 表示；一是切线 $SPT$：$-\infty<h<+\infty\to\{y\text{轴}\}$，其自变量用 $h$ 表示。它的图像一是曲线 $RPQ$，一是直线 $PT$。我们还说现在有两个空间，一是曲线 $RPQ$（或者说它的定义域 $(a,b)$），称为底空间，二是直线 $SPT$（或 $h$ 轴 $\mathbb R^1$）称为切空间。有两个映射 $f:(a,b)\to\mathbb R^1$ 和 $df:\mathbb R_h\to\mathbb R^1$，后者称为前者在 $x$ 点的切映射。魏尔斯特拉斯说的“$h$ 完全不依赖于 $x$”就是说我们研究切映射 $df(x)$ 时，要以 $h$ 为自变量，而 $x$ 则视为固定的，最多也是看作一个参数。我们常把 $x$ 与 $h$ 的变化混在一起分不清楚，因为现在我们讨论的仅是最简单的情况。要研究相应于 $h$ 的函数的全改变量 $f(x+h)-f(x)$ 与 $h$ 之关系，这样就必须而且也能够把这两个不同的变量加起来，甚至给人一个印象即切空间是一个仿射空间：其原点可以任意放到一个位置上（现在是放到底空间的 $x$ 点），而在研究一般的映射与切映射之间的关系时，就不一定（也不需要）能把两个不同的变量加起来，我们研究的也不是 $f(x+h)-f(x)$ 这样的问题。

\begin{figure}[htbp]
\centering
\begin{tikzpicture}[x=0.85cm,y=0.85cm,>=stealth,bella figure]
  \draw[->] (0,0)--(8.5,0) node[right=3pt,fill=none] {$x$};
  \draw[->] (0,0)--(0,5.5) node[above=3pt,fill=none] {$y$};
  \draw[smooth,thick] plot coordinates {(0.8,1.25) (1.5,1.8) (2.5,2.15) (3.6,2.05) (4.5,2.65) (5.4,3.25) (6.35,4.1) (7.1,5.0)};
  \draw[thick] (0.6,1.35)--(7.3,4.8);
  \coordinate (P) at (3.7,2.23);
  \coordinate (R) at (1.35,1.72);
  \coordinate (Q) at (6.9,4.75);
  \coordinate (S) at (0.75,1.42);
  \coordinate (T) at (7.2,4.76);
  \fill (P) circle (1.2pt) node[above left=3pt] {$P$};
  \fill (R) circle (1.2pt) node[left=4pt] {$R$};
  \fill (Q) circle (1.2pt) node[right=4pt] {$Q$};
  \node[below left=2pt] at (S) {$S$};
  \node[above right=2pt] at (T) {$T$};
  \node[above left=2pt] at (2.2,2.18) {$y=f(x)$};
  \draw[dashed] (1.2,0)--(1.2,5.0);
  \draw[dashed] (3.7,0)--(3.7,5.0);
  \draw[dashed] (6.8,0)--(6.8,5.0);
  \node[below=3pt,fill=none] at (1.2,0) {$a$};
  \node[below=3pt,fill=none] at (3.7,0) {$x$};
  \node[below=3pt,fill=none] at (6.8,0) {$b$};
  \draw[->] (0,-0.85)--(8.5,-0.85) node[right=3pt,fill=none] {$h$};
  \node[below=3pt,fill=none] at (3.7,-0.85) {$h=0$};
\end{tikzpicture}
\caption*{图3-2-1}
\end{figure}

以上我们是固定了 $x$ 来看切映射 $df(x)$ 的。如果 $x$ 变动了又如何？例如令 $x$ 变到另一点 $x'$，$P$ 点变成了 $P'$，则切空间（图3-2-2上的虚线）也要变，其原点对着 $x'$ 处。我们把它看成与 $P$ 点的切空间不同的一维空间。这样一来，联系着一小段曲线 $RPQ$，就有许许多多的切空间，各视为互不相干的，这一族切空间称为 $RPQ$ 的切丛。不过我们要注意，一个线性空间画在什么地方是无所谓的，我们也可以把上图中的虚线竖着画（图3-2-2右），使其原点恰好在 $P$ 处，用这个方法来标志它是 $P$ 点的切空间 $T_P$。这样切丛就有了一个形象的表示方法。即令想左图，如果曲线 $RPQ$ 不是平面曲线，则这些“切空间”（即切线）也不会在同一个平面上，因此也不会相交。

\begin{figure}[htbp]
\centering
\begin{tikzpicture}[x=0.75cm,y=0.75cm,>=stealth,bella figure]
  \begin{scope}[xshift=0cm]
    \draw[thick,smooth] plot coordinates {(0.2,0.6) (1.2,1.0) (2.2,1.55) (3.2,2.45) (4.0,3.45)};
    \foreach \x/\y/\s in {0.6/0.78/0.8,1.4/1.1/1.0,2.2/1.55/1.2,3.0/2.25/1.45,3.7/3.1/1.7}{
      \draw[thick] (\x-1.15,\y-0.35*\s)--(\x+1.15,\y+0.35*\s);
    }
    \fill (0.4,0.68) circle (1.1pt) node[left=3pt] {$R$};
    \fill (1.8,1.3) circle (1.1pt) node[below left=3pt] {$P$};
    \fill (2.45,1.73) circle (1.1pt) node[above left=3pt] {$P'$};
    \fill (3.75,3.15) circle (1.1pt) node[right=3pt] {$Q$};
  \end{scope}
  \begin{scope}[xshift=7.0cm]
    \draw[thick,smooth] plot coordinates {(0.0,0.7) (1.0,0.82) (2.0,1.0) (3.0,1.35) (4.0,2.3)};
    \foreach \x in {0.2,0.8,1.4,2.0,2.6,3.2,3.8}{\draw[dashed] (\x,-0.1)--(\x,3.1);}
    \fill (0.25,0.73) circle (1.1pt) node[left=3pt] {$R$};
    \fill (1.45,0.91) circle (1.1pt) node[below left=3pt] {$P$};
    \fill (2.0,1.0) circle (1.1pt) node[above left=3pt] {$P'$};
    \fill (3.8,2.08) circle (1.1pt) node[right=3pt] {$Q$};
  \end{scope}
\end{tikzpicture}
\caption*{图3-2-2}
\end{figure}

切空间、切映射和切丛在现代数学中都有特别重要的地位，我们将在第七章详细讨论。

上面我们提出了一个问题，即 $f'(x)$ 中包含了有关 $f(x)$ 的信息，现在要问，能否从 $f'(x)$ 的信息来恢复 $f(x)$？其实这是很常见的问题：如果已知某一区间 $(a,b)$ 上的 $f'(x)$，则只要再知道 $f(x)$ 在 $(a,b)$ 之某一点 $x_0\in(a,b)$ 处之值 $f(x_0)$，则应有
\[
 f(x)=f(x_0)+\int_{x_0}^x f'(t)dt,\qquad x\in(a,b),
\]
这在物理上是非常自然的：知道了一个质点在某一时刻的位置，又知道它的速度，则此质点的任一时刻的位置自然就完全确定了。然而，我们会在第四章中看见，微分运算的逆运算究竟在什么意义下是积分运算，这决非简单问题。但是现在我们想要问的却是：如果固定了 $x$ 点，则 $f'(x)$ 能在多大程度上给出 $f$ 在此点 $x$ 处的信息？当然，它给出了切线的斜率。但是有不少人常说，若 $f'(x)>0$，则 $f$ 在 $x$ 点单调上升。这句话是很不准确的。因为所谓“上升”、“下降”都是讲的 $f$ 在某一区间 $(a,b)$ 中的动态：若 $x_1,x_2\in(a,b)$，且由 $x_1>x_2$ 可得 $f(x_1)>f(x_2)$，则称 $f$ 在 $(a,b)$ 中上升。说 $f$ 在一点 $x$ 处上升这句话是没有意义的。我们确实可以找到这样的例子，即 $f'(0)>0$，但在包含0的任意小区间中，$f$ 并不上升，$f'$ 且有时正、有时负，从而 $f$ 的切线时而指向上方时而指向下方。与此相类似的还有凸性的概念。例如说 $f''(x)>0$，则 $f$ 在 $x$ 点凸向下是不对的，所谓凸（这里指凸向下）也是关于 $f$ 在某一区间 $(a,b)$ 中的动态：在此区间中任取两点 $a_1,b_1$，如果连接 $(a_1,f(a_1)),(b_1,f(b_1))$ 的弦恒位于连接这两点的弧的上方，则称 $f$ 在 $(a,b)$ 中凸向下。我们这里提醒读者的是，必须注意一个函数在一点处的状态与此函数在一点的某一邻域（不论该邻域如何小）中的状态时常是两回事，读到下面的许多定理时都需要注意于此。

\noindent\textbf{2. 一般情况下微分的定义和性质}\quad 现在考虑 $\mathbb R^m$ 中的开集 $U\subset\mathbb R^m$ 中的映射 $f:U\to\mathbb R^n$。我们限于考虑 $U$ 为开集是为了绕过边界点产生的困难。在1维情况下，开集由最多可数多个开区间 $(a,b)$ 构成。对于一个区间而言，哪怕是闭区间 $[a,b]$，其边界的构造也很简单，无非是两个点。例如对于右端 $b$ 点，我们总可以从它的左方由 $[a,b]$ 内接近于它，而可以考察其左导数。上面我们讲 $f(x)$ 的微分时就遇到过这个问题：若 $x=b$，则必须限于 $h<0$，使 $x+h=b+h$ 仍在 $[a,b]$ 内，但对高维集合 $U$，其边界构造就十分复杂了，甚至无法从 $U$ 的内部去逼近它。所以我们干脆规定 $f$ 只定义在开集 $U$ 上，这时，没有边界点在 $U$ 内，自然就不再有上述困难了。

现在在 $\mathbb R^m$ 中与 $\mathbb R^n$ 中各取坐标系 $x=(x_1,\cdots,x_m),y=(y_1,\cdots,y_n)$，则上述映射可以表示为
\begin{equation}
 y_i=f_i(x_1,\cdots,x_m),\qquad i=1,2,\cdots,n.
 \tag{5}
\end{equation}
特别是 $n=1$ 时，(5)就成了一个 $m$ 元函数
\begin{equation}
 y=f(x_1,\cdots,x_m),
 \tag{6}
\end{equation}
而我们将要讨论的内容就是多元函数的微分学。(5)式中的 $y=(y_1,\cdots,y_n)$ 是一个 $n$ 维向量，所以(5)就成了一个向量值函数，而时常仍用(6)式表示。

现在与定义1平行，我们给出

\noindent\textbf{定义2}\quad 设有两个有限维线性空间 $\mathbb R^m$ 和 $\mathbb R^n$ 以及 $\mathbb R^m$ 之开子集 $U\subset\mathbb R^m$，令 $f:U\to\mathbb R^n$ 是一个（一般为非线性的）映射，若对 $x\in U$，可以找到一个线性映射 $A(x):\mathbb R^m\to\mathbb R^n$，使得
\begin{equation}
 f(x+h)-f(x)=A(x)h+o(h),
 \tag{7}
\end{equation}
则称 $f$ 在 $x$ 点可微。线性映射 $A(x)$ 称为 $f$ 在 $x$ 点的微分，记作 $df(x)$：
\begin{equation}
 A(x)=df(x).
 \tag{8}
\end{equation}
若 $f$ 对 $U$ 之每一点均可微，则说 $f$ 在 $U$ 上可微；若 $A(x)$ 对 $x$ 连续，就说 $f(x)$ 连续可微。

现在举两个容易引起误会的例子。其一是 $f(x)$ 就是线性映射：$f(x)=Ax$。这时
\[
 f(x+h)-f(x)=A(x+h)-Ax=Ah,
\]
而 $o(h)$ 一项变为0。于是 $df=d(Ax)=A$。我们时常就说线性变换（限于常系数）的微分即其自身，但有些文献上线性变换就用 $A$ 表示，$Ax$ 则表示此变换作用于 $x$ 所得之值。这样一来，就可能把这句话误为 $dA=A$，而这个式子是很难解释的。这就涉及第二个例子，设 $f(x)=A$，这是一个常值映射；它把 $x$ 之一切值都映为同一个 $A$。这时 $f(x+h)-f(x)=A-A=0$，而有 $df=dA=0$，这就把上面的 $dA=A$ 弄得更糊涂了。其实，上面一个式子即古典的微积分中的 $d(x)=dx$；下面一个式子则相应于 $dc=0$。这些问题都属于记号问题，或者是把映射和映射作用于某个点所得之值混起来了。下面的定理2后面还有一个类似的说明，希望读者注意。

还要指出，如果 $f$ 在 $x$ 可微，则其微分必是唯一的，因为若有 $A_1,A_2$ 均适合(7)式，必有
\[
 (A_1-A_2)h=o(h).
\]
但是一个由 $\mathbb R^m$ 到 $\mathbb R^n$ 的线性映射除非是0映射，是不可能把任一个 $m$ 维向量 $h$ 都映为高阶无穷小量的（为什么？）所以 $A_1=A_2$，而知 $f$ 之微分若存在必为唯一的。

定义2中说到了 $A(x)$ 对于 $x$ 连续，这句话的意思如下：既然 $A(x):\mathbb R^m\to\mathbb R^n$，则它必可表示为一个 $n$ 行 $m$ 列矩阵 $A=(a_{ij}(x)),i=1,\cdots,n,j=1,\cdots,m$，所谓 $A(x)$ 对 $x$ 连续，现在就是指每一个 $a_{ij}(x)$ 均为 $x$ 在 $U$ 内的连续函数。当然，在更一般的理论框架下，还有进一步的说明。我们知道有限维空间一定可以赋予欧氏空间结构，即对例如 $\mathbb R^m$ 中一点 $x=(x_1,x_2,\cdots,x_m)$ 可以规定其范数
\begin{equation}
 \lVert x\rVert=\left[\sum_{i=1}^m x_i^2\right]^{1/2}.
 \tag{9}
\end{equation}
这样任意两点就有了距离，例如上述 $x$ 点和 $x'=(x'_1,\cdots,x'_m)$ 点的距离就规定为
\[
 \lVert x-x'\rVert=\left[\sum_{i=1}^m(x_i-x_i')^2\right]^{1/2}.
\]
但是，$\mathbb R^m$ 中一点的范数可以用多种方式来定义，例如可以定义为
\begin{equation}
 \lVert x\rVert=\sup_{1\le i\le m}|x_i|.
 \tag{10}
\end{equation}
线性代数理论告诉我们，对于有限维空间，任意两种范数都是等价的，即是说，如果在 $\mathbb R^m$ 中有两个范数 $\lVert\cdot\rVert_1$ 与 $\lVert\cdot\rVert_2$，则必存在两个常数 $c,C>0$ 使对任意一点 $x$ 均有
\[
 c\lVert x\rVert_2\le\lVert x\rVert_1\le C\lVert x\rVert_2.
\]
$c,C$ 与 $x$ 无关。在下面，凡用到 $\lVert\cdot\rVert$ 时，一律是指(9)式所定义的欧几里得范数。这样，(7)式中的高阶无穷小 $o(h)$ 就理解为
\[
 \lVert o(h)\rVert/\lVert h\rVert\to0\qquad(\lVert h\rVert\to0).
\]
下面考虑坐标变换下(7)式如何变化。我们先从线性空间的线性变换开始。$\mathbb R^m$ 和 $\mathbb R^n$ 中有无限多种不同的线性坐标系，如果另取
\[
 x'=(x'_1,\cdots,x'_m)^{\mathrm t}\quad\text{和}\quad y'=(y'_1,\cdots,y'_n)^{\mathrm t},
\]
则必存在两个非奇异的线性变换
\[
 S=(s_{ij})_{m\times m}\quad\text{和}\quad R=(r_{ij})_{n\times n},
\]
使得向量 $x=(x_1,\cdots,x_m)^{\mathrm t}$ 与 $x'$ 以及 $y=(y_1,\cdots,y_n)^{\mathrm t}$ 和 $y'$ 之间有
\[
 {}^{\mathrm t}x'=S\,{}^{\mathrm t}x,\qquad {}^{\mathrm t}y'=R\,{}^{\mathrm t}y.
\]
向量记号左上方的指标 $\mathrm t$ 表示转置，因此例如 $x'$ 等都是竖向量。这时我们有
\[
 {}^{\mathrm t}y'=R\,{}^{\mathrm t}y=R\,{}^{\mathrm t}f(x)=R\,{}^{\mathrm t}f(S^{-1}x')=(R\,{}^{\mathrm t}f\,S^{-1})({}^{\mathrm t}x').
\]
这里我们写 $f$ 当然是表示把 $f$（共有 $n$ 个分量）写成竖向量，其自变量最后写成 $x$ 也是这样。这里没有什么问题。若记
\[
 {}^{\mathrm t}f'=R\,{}^{\mathrm t}f\,S^{-1},
\]
则(7)成 为
\[
 f'(x'+h')-f'(x')=A'h'+o(h').
\]
这里 ${}^{\mathrm t}h'=S\,{}^{\mathrm t}h$ 是一个 $m$ 维向量\jiaokan{原文此处印作“${}^{\mathrm t}h'=S^{-1}{}^{\mathrm t}h$”。但同页前文给出 ${}^{\mathrm t}x'=S{}^{\mathrm t}x$，且下文给出 $A'=RAS^{-1}$；为使坐标变换相容，此处应为 ${}^{\mathrm t}h'=S{}^{\mathrm t}h$（等价地 ${}^{\mathrm t}h=S^{-1}{}^{\mathrm t}h'$）。}，而且与 $\lVert h'\rVert$ 与 $\lVert h\rVert$ 是同阶无穷小，因此 $o(h)$ 与 $o(h')$ 是一样的。$A'=RAS^{-1}$，所以 $A'$ 与 $A$ 是同一个对象在不同坐标系下的表现。通常的微积分教本都只讲到一个多元函数的微分学，或者说值向量 $y$ 只是一个一维向量，这时 $n=1$，从而 $R$ 是一维矩阵，即一个常数，而 $R$ 为非奇异的就意味着此常数不为0。不失一般性，不妨设此常数为1。于是(6)此时成为
\[
 y'=y=f(S^{-1}x')=(fS^{-1})(x'_1,\cdots,x'_m).
\]
这就是一个 $m$ 元函数 $f(x)$ 在坐标变换 $x'=Sx$ 下的结果。既然同一个映射在不同坐标系下有不同的表示，我们应考虑的是它的与坐标无关的，或称内蕴的（coordinate-free, intrinsic）性质。在数学中，寻求一些数学概念或数学量的与坐标无关的性质是一件重要的工作。例如，引进坐标是数学的一大进步，它使得许多数学概念和数学量有了具体的可以操作的形式，甚至可以计算。只要比较一下解析几何与中学里学的初等几何学就可以体会到这一点了。但是，引入坐标系也带来新问题，即把数学量本身的性质与一个特定坐标系的性质混为一谈了。例如看平面上的拉普拉斯算子
\[
 \Delta=\frac{\partial^2}{\partial x^2}+\frac{\partial^2}{\partial y^2}.
\]
如果引入极坐标 $x=r\cos\theta,y=r\sin\theta$，经过适当的计算就会得到
\[
 \Delta=\frac{\partial^2}{\partial r^2}+\frac1r\frac{\partial}{\partial r}+\frac1{r^2}\frac{\partial^2}{\partial\theta^2}.
\]
似乎 $r=0$ 是这个算子的奇点，但是 $r=0$ 相当于直角坐标系中的原点 $(0,0)$，它并不是 $\Delta=\frac{\partial^2}{\partial x^2}+\frac{\partial^2}{\partial y^2}$ 的奇点。产生这个现象的原因是：原点本身就是极坐标的奇点；它的极坐标并非唯一的，而是任意的 $(0,\theta)$，$\theta$ 可以取任意值。

关于映射的微分的坐标无关的讨论见下一段。现在我们仍先固定 $\mathbb R^m$ 和 $\mathbb R^n$ 的坐标系 $x$ 与 $y$，并对这个坐标系讨论映射 $f$ 可微的必要充分条件，并求出 $A(x)$ 的表达式。

\noindent\textbf{定理1}\quad 定义2中的映射 $f$ 在某坐标系下在开集 $\Omega$ 中连续可微的必要充分条件是，$f$ 的各个分量 $y_i$ 对任意 $x_j$ 均为连续可微函数。这时
\begin{equation}
 A(x)=\frac{\partial(y_1,\cdots,y_n)}{\partial(x_1,\cdots,x_m)}=
 \begin{pmatrix}
 \dfrac{\partial y_1}{\partial x_1}&\cdots&\dfrac{\partial y_1}{\partial x_m}\\
 \vdots&&\vdots\\
 \dfrac{\partial y_n}{\partial x_1}&\cdots&\dfrac{\partial y_n}{\partial x_m}
 \end{pmatrix}.
 \tag{11}
\end{equation}
(11)右方矩阵称为雅可比矩阵（Jacobian matrix）。

\noindent\textbf{证}\quad 先证充分性，我们不妨只看 $n=1,m=2$ 的特例，这时 $h=(h_1,h_2)$ 而且映射 $f$ 就是 $f=f_1$，于是设它是一个连续可微函数：
\begin{align}
 f(x+h)-f(x)&=f(x_1+h_1,x_2+h_2)-f(x_1,x_2)\\
 &=\bigl[f(x_1+h_1,x_2+h_2)-f(x_1,x_2+h_2)\bigr]\\
 &\quad+\bigl[f(x_1,x_2+h_2)-f(x_1,x_2)\bigr].
 \tag{12}
\end{align}
当 $f$ 对 $(x_1,x_2)$ 在开集 $\Omega$ 中均为连续可微时，对右方两项分别应用拉格朗日公式有
\[
 f(x+h)-f(x)=\frac{\partial f}{\partial x_1}(x_1+\theta_1h_1,x_2+h_2)h_1+\frac{\partial f}{\partial x_2}(x_1,x_2+\theta_2h_2)h_2.
\]
再利用 $\dfrac{\partial f}{\partial x_1},\dfrac{\partial f}{\partial x_2}$ 在 $(x_1,x_2)$ 处的连续性即得
\[
 f(x+h)-f(x)=\left(\frac{\partial f}{\partial x_1}(x)h_1+\frac{\partial f}{\partial x_2}(x)h_2\right)+o(h)
 =A(x)h+o(h).
\]
而这里 $A(x)$ 是 $1\times2$ 矩阵 $\left(\dfrac{\partial f}{\partial x_1},\dfrac{\partial f}{\partial x_2}\right)$。由假设，它是连续的，于是映射 $f$ 在此开集中连续可微，而且其微分就是雅可比矩阵，充分性得证。

必要性很简单。设 $f$ 在此开集中连续可微，任取此开集中一点 $x$，对 $f$ 之任一分量（不妨即设为 $f_1=y_1$），对 $h=(h_1,0)$ 有
\[
 f(x_1+h_1,x_2)-f(x_1,x_2)=A(x)h_1+o(h).
\]
双方除以 $h_1$ 再令 $h_1\to0$，即知 $\dfrac{\partial f}{\partial x_1}$ 连续。同理 $f$ 之各分量对各个 $x_i$ 之偏导数 $\dfrac{\partial y_j}{\partial x_i}$ 均在此开集中连续而必要性得证。

\noindent\textbf{注1}\quad 在 $U$ 中连续可微的映射 $f$ 之集合构成一个空间。我们通常记作 $C^1(U)$。上面说到了这时 $df=A(x)$ 是雅可比矩阵，而且其元素都是古典意义下的连续可微函数。这类函数之集合也记作 $C^1(U)$。这样看来，$C^1(U)$ 的两种说法是一致的，我们不必加以区别。不过，后一种意义下说 $A(x)$ 之元素是 $C^1$ 函数是在选定了一个坐标系以后说的，如果变一个坐标系，它是否仍在 $C^1(U)$ 中？答案是肯定的，下面会详细解释。

\noindent\textbf{注2}\quad 定义2中讲到了 $f$ 在一点 $x$ 可微，因此有一个错觉，以为其必要充分条件应是 $\dfrac{\partial y_i}{\partial x_j}$ 在 $x$ 点存在。实际上不是这样，而只当 $m=1$ 时它是对的，因为 $m=1$ 时，1维向量 $h$ 就是一个数 $h$，向量 $h\ne0$ 就是数 $h\ne0$。于是对在定义1的(4)式（那里 $n$ 也是1，实际上 $n>1$ 也对）双方可以用 $h$ 去除，而当 $h\to0$ 时立即有 $A=\dfrac{df(x)}{dx}=f'(x)$。$m>1$ 时我们就不能用 $m$ 维向量 $h$ 去除(7)式两侧，而要利用(12)式过渡到一个自变量的情况，并且应用拉格朗日公式。这就要用到 $x$ 点以外其它点处偏导数的存在，然后还要利用其在 $x$ 点的连续性，把这些偏导数化为在 $x$ 点的值再加上一个无穷小量。这就看到，$m>1$ 时可微性与可求导数是不相同的，二者有本质区别。但是我们无法把可微性准确地刻画为可求导，再加上导数应该合什么条件，所以，我们干脆讨论某一开集中函数的连续可微性，而把我们需要的结果以定理1的形式陈述出来。这样做不但足够用了，而且还想借此再次提醒读者：一个函数（即一个映射）在一点的情况与在一点的某个邻域中的情况是不一样的。

于是若我们确定了一个坐标系 $(x_1,\cdots,x_m)$，则点 $x$ 的变化可以是沿某一个坐标轴，例如 $x_1$ 轴。所以 $h=(h_1,0,\cdots,0)$，而我们可以写出
\begin{equation}
 \Delta_1 f=f(x_1+h_1,x_2,\cdots,x_m)-f(x)=A_1(x)h_1+o(h_1).
 \tag{13}
\end{equation}
这样的 $A_1(x)$ 相应于偏导数，或称为偏微分。为了 $f$ 有偏微分，由极限式
\[
 \lim_{h_1\to0}\frac{\Delta_1f}{\Delta x_1}=A_1(x),\qquad \Delta x_1=h_1,
\]
可知只要 $f$ 为可求导即可。与偏微分相对立，(7)中的 $A(x)$ 称为全微分。偏微分亦称加托（Gateaux）微分，全微分则称弗雷歇（Fréchet）微分。二者的关系我们就不讲了。但由此容易理解为什么 $m=1$ 与 $m>1$ 时可微性的情况不一样。

\noindent\textbf{注3}\quad 现在，微分是一个矩阵，线性代数告诉我们一个 $n\times m$ 矩阵 $A$，就是一个 $\mathbb R^m$ 到 $\mathbb R^n$ 的线性映射。我们把映射 $f:\Omega\to\mathbb R^n$ 的 $df$ 微分也定义为一个映射。当 $m=n=1$ 时，我们会有一个一阶矩阵，但是一阶矩阵也就是一个数，所以这时 $A(x)=(df)(x)$ 按古典的微分的观点看来是一个数，即导数 $f'(x)$。但是按我们现在的观点看来则是一个映射 $df:\mathbb R^1\to\mathbb R^1$。按古典的观点 $df=f'(x)dx=f'h$，按现在的观点则 $df=f'(x)$，所以在一些文献中 $df$ 就表示 $f'(x)$。这种记号上的混淆只有请读者自己注意。倒是值得注意的是，$df$ 在固定的坐标系下要用雅可比矩阵
\begin{equation}
 J(x)=\frac{\partial(y_1,\cdots,y_n)}{\partial(x_1,\cdots,x_m)}=
 \begin{pmatrix}
 \dfrac{\partial y_1}{\partial x_1}&\dfrac{\partial y_1}{\partial x_2}&\cdots&\dfrac{\partial y_1}{\partial x_m}\\
 \dfrac{\partial y_2}{\partial x_1}&\dfrac{\partial y_2}{\partial x_2}&\cdots&\dfrac{\partial y_2}{\partial x_m}\\
 \vdots&\vdots&&\vdots\\
 \dfrac{\partial y_n}{\partial x_1}&\dfrac{\partial y_n}{\partial x_2}&\cdots&\dfrac{\partial y_n}{\partial x_m}
 \end{pmatrix}
 \tag{14}
\end{equation}
表示，于是也就可以说雅可比矩阵是 $f'(x)$ 的合适的推广。

还有一点要提醒，在古典的微积分教本中是在 $m=n$ 时讲雅可比行列式，并且采用下面的记号
\begin{equation}
 J(x)=\frac{\partial(y_1,\cdots,y_n)}{\partial(x_1,\cdots,x_n)}=
 \begin{vmatrix}
 \dfrac{\partial y_1}{\partial x_1}&\cdots&\dfrac{\partial y_1}{\partial x_n}\\
 \vdots&&\vdots\\
 \dfrac{\partial y_n}{\partial x_1}&\cdots&\dfrac{\partial y_n}{\partial x_n}
 \end{vmatrix},
 \tag{15}
\end{equation}
但 $m\ne n$ 时就谈不到雅可比行列式，所以我们只讲雅可比矩阵，$J(x)$ 与 $\frac{\partial(y_1,\cdots,y_n)}{\partial(x_1,\cdots,x_m)}$ 则按(14)理解，而把雅可比行列式写作 $\det J(x)$ 或 $\det\frac{\partial(y_1,\cdots,y_n)}{\partial(x_1,\cdots,x_m)}$，这些也要读者多加小心。

关于为什么雅可比矩阵是 $f'(x)$ 的合适的推广，下面还要详细讲。
